\section{変分量子アルゴリズム}
\label{sec:variational-quantum-algorithm}

\subsection{深い量子回路とNISQ}
\label{subsec:deep-circuits-and-nisq}

量子誤り訂正を前提とする大規模な量子アルゴリズムでは，多数の量子ビットと深い量子回路が必要になる．
一方，現在利用可能な量子デバイスは量子ビット数，ゲート精度，コヒーレンス時間などに制約があり，
ノイズの影響を受けやすい．
このような中小規模で誤り訂正を行わない量子デバイスは，
NISQ（Noisy Intermediate-Scale Quantum）デバイスと呼ばれる．

NISQデバイスでは，量子回路を短く保ちつつ量子計算を利用する方法が重要になる．
その代表的な考え方が，量子回路による状態生成と測定を古典計算による最適化と組み合わせる
変分量子アルゴリズムである．

\subsection{古典・量子ハイブリッド計算}
\label{subsec:hybrid-computation}

古典・量子ハイブリッド計算では，量子デバイスと古典計算機が次の処理を繰り返す．

\begin{enumerate}[(1)]
  \item 古典計算機が回路パラメータを設定する．
  \item 量子デバイスがパラメータ付き量子回路を実行する．
  \item 測定結果または観測量の期待値を古典計算機へ返す．
  \item 古典計算機が目的関数を評価し，回路パラメータを更新する．
  \item あらかじめ定めた停止条件を満たすまで処理を繰り返す．
\end{enumerate}

量子デバイスは高次元の量子状態を生成して測定し，古典計算機は損失関数の評価と
パラメータ更新を担当する．
この役割分担によって，深い量子回路を避けながら量子回路を学習モデルの一部として利用できる．

\subsection{パラメータ付き量子回路}
\label{subsec:parameterized-quantum-circuit}

入力データ$x$を符号化する回路を$U_{\mathrm{in}}(x)$，
学習対象のパラメータ$\bm{\theta}$を持つ量子回路を$U(\bm{\theta})$とする．
初期状態を$\ket{0}^{\otimes n}$とすると，回路の出力状態は
\begin{equation}
  \ket{\psi(x,\bm{\theta})}
  =
  U(\bm{\theta})
  U_{\mathrm{in}}(x)
  \ket{0}^{\otimes n}
  \label{eq:parameterized-quantum-state}
\end{equation}
と表される．
観測量$M$の期待値をモデルの出力とすれば，
\begin{equation}
  f(x,\bm{\theta})
  =
  \bra{\psi(x,\bm{\theta})}
  M
  \ket{\psi(x,\bm{\theta})}
  \label{eq:pqc-output}
\end{equation}
となる．

式\eqref{eq:pqc-output}は，量子回路学習で用いられる基本的なモデルの形である．
入力$x$は量子状態へ符号化され，$\bm{\theta}$は学習によって更新される．
次章では，教師あり学習の枠組み，損失関数，パラメータ更新方法を導入し，
このパラメータ付き量子回路を回帰モデルとして利用する方法を説明する．

